\documentclass[a4paper,9pt]{extarticle}

%=========================================================
% PACKAGES
%=========================================================

\usepackage[
a4paper,
top=1.8cm,
bottom=1.5cm,
footskip=7mm,
left=1cm,
right=1cm,
headheight=14pt,
headsep=5mm
]{geometry}
\usepackage[T1]{fontenc}
\usepackage[utf8]{inputenc}
\usepackage[french]{babel}

\usepackage[sfdefault]{FiraSans}
\renewcommand{\familydefault}{\sfdefault}

\usepackage{multicol}
\setlength{\columnsep}{0.7cm}

\usepackage{amsmath,amssymb}

\usepackage{titlesec}
\usepackage{enumitem}

\usepackage{tcolorbox}
\tcbuselibrary{skins,breakable}

\usepackage{xcolor}

\usepackage{graphicx}
\usepackage{float}

\usepackage{tikz}

\usepackage{fancyhdr}
\usepackage[hidelinks]{hyperref}


%=========================================================
% PAGE SETTINGS
%=========================================================

\setlength{\parindent}{0pt}
\setlength{\parskip}{1.5pt}

\setlength{\abovedisplayskip}{4pt}
\setlength{\belowdisplayskip}{4pt}

\setlist[itemize]{
leftmargin=*,
itemsep=1pt,
topsep=1pt
}


%=========================================================
% COLORS
%=========================================================

\definecolor{Blue}{HTML}{E3F0FF}
\definecolor{Green}{HTML}{E3F7EA}
\definecolor{Pink}{HTML}{FFE5ED}
\definecolor{Orange}{HTML}{FFF0DC}
\definecolor{Purple}{HTML}{EEE7FF}
\definecolor{Yellow}{HTML}{FFFADB}


\definecolor{BlueTitle}{HTML}{285B8F}
\definecolor{GreenTitle}{HTML}{287052}
\definecolor{PinkTitle}{HTML}{A84063}
\definecolor{OrangeTitle}{HTML}{A96500}
\definecolor{PurpleTitle}{HTML}{6245A5}


%=========================================================
% HEADER
%=========================================================

\pagestyle{fancy}

\fancyhf{}

\fancyhead[L]{\small\textbf{Computer Security}}
\fancyhead[C]{\small Access Control}
\fancyhead[R]{\small Week 03 --- 21-27/09/2026}

\fancyfoot[C]{\small\href{https://github.com/AlphaXT38G/notes_ba5}{github.com/AlphaXT38G/notes\_ba5}}
\fancyfoot[R]{\small \thepage}

\renewcommand{\headrulewidth}{0.4pt}



%=========================================================
% TITLES
%=========================================================

\titleformat{\section}
{\large\bfseries\color{BlueTitle}}
{}{0pt}{}

\titlespacing{\section}
{0pt}{6pt}{3pt}


% separator

\newcommand{\separator}{
\vspace{2pt}
\hrule
\vspace{3pt}
}


%=========================================================
% BOXES
%=========================================================

% DEFINITION - blue
\newtcolorbox{definition}{
enhanced,
breakable,
colback=Blue,
colframe=BlueTitle,
boxrule=0.7pt,
arc=0mm,
title=\textbf{Definition},
fonttitle=\bfseries
}

% IMPORTANT - pink
\newtcolorbox{important}{
enhanced,
breakable,
colback=Pink,
colframe=PinkTitle,
boxrule=0.7pt,
arc=0mm,
title=\textbf{Key idea},
fonttitle=\bfseries
}

% FORMULAS - green
\newtcolorbox{formula}{
enhanced,
breakable,
colback=Green,
colframe=GreenTitle,
boxrule=0.7pt,
arc=0mm,
title=\textbf{Key formulas},
fonttitle=\bfseries
}

% WARNING - orange
\newtcolorbox{warning}{
enhanced,
breakable,
colback=Orange,
colframe=OrangeTitle,
boxrule=0.7pt,
arc=0mm,
title=\textbf{Warning},
fonttitle=\bfseries
}

% EXAMPLE - purple
\newtcolorbox{example}{
enhanced,
breakable,
colback=Purple,
colframe=PurpleTitle,
boxrule=0.7pt,
arc=0mm,
title=\textbf{Example},
fonttitle=\bfseries
}

% SUMMARY - orange
\newtcolorbox{summary}{
enhanced,
breakable,
colback=Yellow,
colframe=OrangeTitle,
boxrule=0.7pt,
arc=0mm,
title=\textbf{Summary},
fonttitle=\bfseries
}



%=========================================================
% FIGURE BOX
%=========================================================

\newtcolorbox{figurebox}{
enhanced,
colback=white,
colframe=gray,
boxrule=0.5pt,
arc=0mm,
center
}



%=========================================================
% DOCUMENT
%=========================================================

\begin{document}

%\footnotesize   --> To reduce the size of the text

\thispagestyle{fancy}



\begin{multicols}{2}


%=========================================================
\section{Access Control}


\begin{definition}

\textbf{ACCESS CONTROL:} Security mechanism that ensures that
``all accesses and actions on objects by principals are WITHIN the security policy''

\end{definition}


\textbf{DISCRETIONARY ACCESS CONTROL (DAC)}

\begin{itemize}

\item Object owners assign permissions

\item Ownership of resources

\item Windows, Linux, macOS

\item Social Networks

\end{itemize}


\separator


\textbf{MANDATORY ACCESS CONTROL (MAC)}

Central security policy assigns permissions

\begin{itemize}

\item Organizations with need for central controls

\item Military -- focus on confidentiality

\item Hospital environment -- focus on confidentiality
and integrity

\item Banking -- focus on integrity

\end{itemize}



%=========================================================
\section{DAC}


\begin{definition}

\textbf{ACCESS CONTROL MATRIX:} an abstract representation of all permitted triplets
of (subject, object, access right) within a system

\end{definition}


\begin{itemize}

\item \textbf{Subjects (principals):} entity within an IT system

a user, a process, a service

\item \textbf{Objects (assets):} resources that (some) subject may access or use

a file, a folder, a row in a database, the system's memory, a machine in the network,
a printer, a page in a website

\item \textbf{Operations:} in abstract, subjects can observe and/or alter objects

read, write, append, execute

\end{itemize}



%=========================================================
\section{Role Based Access Control (RBAC)}


\begin{important}

Subjects are similar to each other: assign same rights

\end{important}


\begin{example}

e.g., a doctor has the same privileges as another doctor

\end{example}


\begin{enumerate}[leftmargin=*, itemsep=1pt, topsep=1pt]

\item assign permissions to roles

\item assign roles to subjects

\item subjects select an active role -- they have the permissions of the active role

\end{enumerate}


\begin{warning}

\textbf{RBAC Problems}

\textbf{Problem 1: Role Explosion}

\begin{itemize}

\item Temptation to create fine grained roles, denying benefits of RBAC

\end{itemize}

\textbf{Problem 2: Simple RBAC has limited expressiveness}

\begin{itemize}

\item Problems with implementing least privilege

\item Some roles are relative: ``Alice's Doctor'' vs. ``Any Doctor''

\end{itemize}

\textbf{Problem 3: Difficult to implement separation of duty}

\begin{itemize}

\item ``Two doctors are needed to authorize a procedure''

\item RBAC Mechanism needs to ensure they are distinct!

\end{itemize}

\end{warning}



%=========================================================
\section{Group Based Access Control}


\begin{important}

\textbf{Observation:} Some permissions are always needed together

\end{important}


\begin{example}

e.g., access to sockets and network interface always go hand in hand

\end{example}


\begin{enumerate}[leftmargin=*, itemsep=1pt, topsep=1pt]

\item assign permissions to access objects to groups

\item assign subjects to groups

\item subjects have the permissions of all their groups

\end{enumerate}



%=========================================================
\section{Ambient Authority}


\begin{definition}

\textbf{AMBIENT AUTHORITY} is used by a subject if for an action to succeed it only needs to
specify the operation and the names of the involved object(s)

\end{definition}


\begin{itemize}

\item no need to repeat the subject all the time (usability)

\item least privilege becomes harder to enforce

\end{itemize}


\begin{warning}

confused deputy problem! PK ????

\end{warning}


\separator


\textbf{The confused deputy problem}

Problem with ambient authority:

\begin{important}

A privileged program can be tricked to misuse its authority

(Confused deputy problem)

\end{important}


\textbf{Solutions:}

\begin{enumerate}[leftmargin=*, itemsep=1pt, topsep=1pt]

\item Re-implement access control in the privileged process

\item Let privileged process check authorization for Alice.

\item Capabilities can help!

\end{enumerate}


\separator


\textbf{Files have ACLs attached to them}

\begin{itemize}

\item Each file is assigned an owner UID and GID $\rightarrow$ GID???

\item Each file has 9 permission bits

\end{itemize}



\end{multicols}


\end{document}
